\section{Block-theoretic bounds and transfer}
\label{sec:preliminaries}

We first isolate the block-theoretic transfers used in the two
parts of the proof.  The ordinary-coordinate estimate recovers full
projective-character bounds from selected rows.  The central,
restriction and abelian-extension estimates then carry those
bounds between the groups in the quasisimple reduction.

Fix a prime \(p\), put \(k=\overline{\F}_p\), and let
\(\cO=W(k)\) be its Witt ring.  The letter \(M\) denotes an upper bound
on defect-group order.  Unless a further parameter is displayed, a
uniform bound may depend on \(p\) and \(M\), but not on the ambient
finite group.  When ordinary characters are used, we extend scalars to
a splitting characteristic-zero field.  This leaves the decomposition numbers unchanged and extends each
chosen integral Morita equivalence.  We do not infer descent of an
arbitrary equivalence from a splitting extension.

For a finite group \(G\), a block is a nonzero primitive central
idempotent \(b\) of \(kG\), or the corresponding summand \(kGb\).
The idempotent lifts uniquely to a block of \(\cO G\), denoted by the
same letter.  A defect group is a maximal \(p\)-subgroup \(D\) for which
\(\Br_D(b)\ne0\).  Here
\[
 \Br_D:(kG)^D\longrightarrow kC_G(D)
\]
deletes coefficients of group elements outside \(C_G(D)\).
Morita equivalences throughout are linear over the indicated
coefficient ring.

Write \(\Irr(b)\) for the ordinary irreducible characters in \(b\),
\(\IBr(b)\) for its irreducible Brauer characters, and
\(k(b)=|\Irr(b)|\), \(l(b)=|\IBr(b)|\).  If
\(\chi^0\) is the restriction of \(\chi\) to \(p\)-regular elements, then
\[
 \chi^0=\sum_{\varphi\in\IBr(b)}d_{\chi\varphi}\varphi .
\]
The nonnegative integral matrix
\(D_b=(d_{\chi\varphi})\) is the decomposition matrix.
Brauer reciprocity identifies its columns with the ordinary
characters \(\Phi_\varphi\) of projective indecomposable modules
(PIMs).  Thus the Cartan matrix is
\[
 C_b=D_b^{\mathsf T}D_b,\qquad
 (C_b)_{\varphi\psi}=\langle\Phi_\varphi,\Phi_\psi\rangle_G.
\]
In particular, a bound on PIM norms bounds every Cartan entry by
Cauchy--Schwarz.

We use the following standard facts of block theory
\cite{Linckelmann,BrauerFeit}.  The Brauer--Feit theorem bounds \(k(b)\)
in terms of defect-group order, and \(l(b)\le k(b)\).  The elementary
divisors of \(C_b\) divide \(|D|\); in particular
\[
 0<\det C_b\le |D|^{l(b)}.
\]
There is an ordinary character of height zero in every block.
If \(N\lhd G\) and \(B\) covers a block \(b\) of \(N\), defect groups
can be chosen so that \(D_b=D_B\cap N\).  Fong--Reynolds reduction
replaces the covering problem by the inertia group of \(b\), preserving
Morita type and defect.  If \(b\) is invariant and \(G/N\) is a
\(p\)-group, then \(b\) remains a block of \(G\) and
\[
 D_BN=G.
\]
These assertions will only be used in these stated forms.  We also
use Clifford theory for restriction to a normal subgroup, including
its projective multiplicity-space formulation.

\subsection{Central quotients and ordinary coordinates}

\begin{lemma}[Central transfer]
\label{lem:central-transfer}
Let \(Z\le Z(G)\) be a \(p\)-subgroup.  Blocks of \(G\) correspond to
blocks of \(G/Z\), every defect group contains \(Z\), and corresponding
defect orders differ by \(|Z|\).  Their simple modules correspond by
inflation.  With these identifications their Cartan matrices satisfy
\[
 C_b=|Z|\,C_{\bar b}.
\]
In particular, Cartan bounds transfer in either required direction
when \(|Z|\) is bounded.  For a central \(p'\)-subgroup \(Z'\), the
summand on which \(Z'\) acts trivially is the averaging-idempotent
summand and identifies with the group algebra of \(G/Z'\).
\end{lemma}

\begin{proof}
The block and defect assertions are the central-quotient theorem.
For the Cartan assertion put \(J=\rad(kZ)\).
The ideal \(JkG\) is nilpotent and its quotient is \(k(G/Z)\), so
primitive projective covers correspond under passage to coinvariants.
If \(P\) is such a cover, then \(P\) is projective, hence free, as a
\(kZ\)-module.  The canonical multiplication maps give
\[
 (J^i/J^{i+1})\otimes_k(P/JP)
       \ \simeq\ J^iP/J^{i+1}P .
\]
They are isomorphisms of \(G/Z\)-modules: centrality makes the action
on the first factor trivial, and changing a lift in \(G\) changes
the action only by the next radical layer.  The sum of the dimensions
of \(J^i/J^{i+1}\) is \(|Z|\).  Additivity of composition multiplicities
therefore gives the displayed equality.  For \(Z'\), use
\(|Z'|^{-1}\sum_{z\in Z'}z\).
\end{proof}

A set \(\mathcal B\subseteq\Irr(b)\) is an \emph{ordinary integral
basic set} if the characters \(\{\chi^0:\chi\in\mathcal B\}\)
are a \(\Z\)-basis of the Brauer character lattice.
The corresponding square submatrix of \(D_b\) then has determinant
\(\pm1\).  We also use basic sets for a direct sum of blocks.

\begin{lemma}[Bounded inverse projection]
\label{lem:row-projection}
For blocks of defect order at most \(M\), including unions of a
bounded number of such blocks, let \(V\) be the rational span of
their PIM columns in ordinary-character coordinates.
If a coordinate projection \(\pi\) is injective on \(V\), its inverse
\[
 (\pi|_V)^{-1}:\pi(V)\longrightarrow V
\]
has uniformly bounded Euclidean operator norm.  All full-column-size
minors of a decomposition matrix, before or after coordinate
projection, are uniformly bounded.

If a basic set for a union of blocks has bounded size and the
\(p\)-defects of all its ordinary degrees are bounded, then the number
and defect orders of those blocks, and hence the dimensions of their full ordinary-character coordinate
spaces, are bounded.  This bounds the numbers of ordinary irreducible
characters, not their degrees.
\end{lemma}

\begin{proof}
Let the PIM columns be \(v_1,\ldots,v_l\in\Z^n\).  Their exterior
product is a nonzero integral vector, and Cauchy--Binet gives
\[
 \|v_1\wedge\cdots\wedge v_l\|^2
    =\det(D_b^{\mathsf T}D_b)
    =\sum_{|I|=l}\det(D_{b,I})^2 .
\]
The preceding standard bounds control \(n,l\) and this determinant.
After padding the ambient coordinate spaces to one fixed dimension,
only finitely many exterior products occur.  Each determines the
subspace \(V\).  There are also only finitely many coordinate
projections in that dimension.  Restricting to the injective ones,
the maximum of their inverse norms is finite.  The same identity
bounds each displayed minor.  For a bounded block union take the
orthogonal direct sum.

A basic set partitions by block.  Indeed each ordinary character
has Brauer constituents in its own block, and the Brauer lattice is
the direct sum of the block lattices.  The basis property consequently
restricts to each summand, which must have at least one basic row.
It remains to recover its defect from those rows.
Fix a block \(b\) in the union and put
\(\mathcal B_b=\mathcal B\cap\Irr(b)\).
Write \(|G|_p=p^a\), and let \(b\) have defect \(p^d\).
Every ordinary degree in the block has \(p\)-valuation at least
\(a-d\).  Some degree has valuation exactly \(a-d\), by height zero.
If every basic-row degree had larger valuation, evaluating its
integral spanning expression at \(1\) would give the same larger
valuation for every ordinary degree, a contradiction.  Thus
\[
 d=\max_{\chi\in\mathcal B_b}
           v_p\bigl(|G|/\chi(1)\bigr).
\]
The asserted bounds now follow from the bound on basic-set size
and the Brauer--Feit theorem.
\end{proof}

This lemma bounds the projective \emph{subspace}, not a selected
nonnegative integral basis of it.  That distinction lets us first
estimate selected ordinary coordinates and only afterwards recover
the full PIM norms.

\subsection{Crossed algebras and upward Cartan transfer}

We record conventions for crossed algebras that will also be used in
Section~\ref{sec:extensions}.  A crossed algebra on a finite-dimensional
\(k\)-algebra \(R\), graded by a finite group \(A\), has the form
\[
 T=\bigoplus_{x\in A}Ru_x,\qquad
 u_xr=\alpha_x(r)u_x,\qquad u_xu_y=a_{xy}u_{xy},
\]
where each \(u_x\) is invertible, \(\alpha_x\in\Aut(R)\),
and \(a_{xy}\in R^\times\).  Associativity means
\[
 \alpha_x\alpha_y=\operatorname{ad}(a_{xy})\alpha_{xy},\qquad
 a_{xy}a_{xy,z}=\alpha_x(a_{yz})a_{x,yz}.
\]
In the second equality \(a_{xy,z}\) denotes the factor with first
index the group product \(xy\).  Replacing \(u_x\) by \(v_xu_x\),
\(v_x\in R^\times\), is called a change of homogeneous units.

If \(e\) is a full basic idempotent of \(R\), then
\(\alpha_x(e)\) is unit-conjugate to \(e\): both select one copy of
each indecomposable projective type.  Adjusting \(u_x\) by such units
makes it commute with \(e\).  Thus \(eTe\) is again crossed on \(eRe\),
and \(e\) is full in \(T\).  In particular
\[
 \dim_k(eTe)=|A|\dim_k(eRe).
\]
The same argument applies to block orders and integral basic
idempotents.

\begin{lemma}[Transfer across an abelian quotient]
\label{lem:abelian-transfer}
Suppose \(N\lhd G\) and \(G/N\) is abelian.  Among blocks \(B\) of
\(G\) of defect order at most \(M\), uniform Cartan bounds for their
covered blocks of \(N\) imply uniform Cartan bounds for \(B\).
No bound on the \(p'\)-part of \(|G/N|\) is required.
\end{lemma}

\begin{proof}
Apply Fong--Reynolds and assume the covered block \(b\) is invariant.
Let \(G_p/N\) be the Sylow \(p\)-subgroup of \(G/N\).  The unique
block of \(G_p\) above \(b\) has a defect group surjecting onto
\(G_p/N\); its order is bounded by a defect group of \(B\).
Consequently \(|G_p/N|\le M\).  A basic corner of the crossed
extension by this bounded group has bounded dimension, by the
preceding observation.  Its Cartan entries are therefore bounded.
We may replace \(N\) by \(G_p\), leaving a \(p'\)-quotient \(A\).
The identity block is invariant; alternatively one may take its
inertia group again.

Pass to a full basic corner, and write \(R\) for the identity
algebra and \(T\) for the resulting crossed algebra.
The dimension of \(R\) is bounded: for a basic algebra it is the
sum of its Cartan entries, and its number \(n\) of simple modules
is bounded by defect.  Put
\[
 d=\max_{i,j}\dim_k\Ext^1_R(S_i,S_j),\qquad
 (\rad R)^L=0.
\]
Both \(d\) and \(L\) are bounded.  Since \(\rad(R)\) is invariant,
\(\rad(R)T\) is nilpotent.  Maschke's theorem for the crossed
algebra on the semisimple quotient gives
\[
 \rad(T)=\rad(R)T,\qquad (\rad T)^L=0.
\]

The radical length is now controlled.  To bound the basic algebra
of the chosen block, it remains to bound the numbers of arrows in
its quiver, that is, first extensions between its simple modules.
Let \(X,Y\) be simple modules in the block of \(T\) under
consideration.  Their restrictions are semisimple over \(R\):
the nilpotent ideal \(\rad(R)T\) annihilates both modules.
Clifford theory describes the support of each restriction as one
\(A\)-orbit of simple \(R\)-types.  At each type its multiplicity
space is an irreducible projective representation of the stabilizer.
The scalar factor system accounts for the chosen intertwiners and
the inner factors \(a_{xy}\).  Simplicity makes that projective
representation irreducible.

A \(T\)-projective resolution restricts to an \(R\)-projective
resolution.  The group action on its \(R\)-linear Hom complex is
genuine, since the inner crossed factors cancel by \(R\)-linearity.
As \(|A|\) is invertible in \(k\), taking invariants is exact.
Thus \(\Ext^1_T(X,Y)\) is computed by invariant parts of the
extension spaces between these semisimple restrictions.
Decompose those spaces into orbits of pairs of simple types.
For a representative pair \(i,j\), write the corresponding
isotypic summands of \(X,Y\) as \(S_i\otimes V\) and
\(S_j\otimes W\), respectively.  Put \(H=A_i\cap A_j\) and
\(E=\Ext^1_R(S_i,S_j)\), with the induced projective
\(H\)-action.  Restrict \(V,W\) from their stabilizers to \(H\).
The corresponding summand has dimension
\[
 \dim\Hom_H(V,W\otimes E),
 \qquad \dim E\le d,
\]
with the compatible projective factors understood.

For completeness, these multiplicity spaces need not have bounded
dimensions.  Instead their normalized character norms give
\begin{equation}
 \dim\Hom_H(V,W\otimes E)
 \le \dim(E)\sqrt{[A_i:H][A_j:H]} .
 \label{eq:ext-stabilizer-bound}
\end{equation}
To see this, normalize the scalar factors to roots of unity of
\(p'\)-order and realize the projective representations on compatible
finite central \(p'\)-extensions.  Such normalization is possible
because scalar cohomology is killed by the group order.  These
semisimple representations lift to characteristic zero.
On the common scalar-character summands, the character formula,
\(|\chi_E(h)|\le\dim E\), and Cauchy--Schwarz reduce the estimate to
\[
 \frac1{|H|}\sum_{h\in H}|\chi_V(h)|^2\le[A_i:H],
 \qquad
 \frac1{|H|}\sum_{h\in H}|\chi_W(h)|^2\le[A_j:H].
\]
These follow by enlarging the sums to the respective stabilizers,
where the normalized squared norms are \(1\).
This proves \eqref{eq:ext-stabilizer-bound}.

The two indices are at most \(n\), since they count orbits of
simple types under a subgroup of \(A\); there are at most \(n^2\)
pair-orbit contributions.  Hence the coarse uniform estimate
\[
 \dim\Ext^1_T(X,Y)\le dn^3
\]
suffices.  The block of \(T\) corresponding to \(B\) has a bounded
number \(m\) of simple modules by its defect bound.
Its basic algebra has \(m\) vertices and at most \(m^2dn^3\)
arrows.  Lifts of a basis of its radical modulo its square generate
the radical, so paths of length less than \(L\) span that algebra.
Its dimension, and therefore its Cartan entries, are bounded.
\end{proof}

\subsection{Restriction and decomposition rows}

The preceding lemma transfers Cartan bounds upwards even across
an abelian quotient of unbounded $p'$-order.  The next lemma supplies
the complementary downward estimate: it controls decomposition
numbers through restriction multiplicities, without requiring a
bound on the index itself.

\begin{lemma}[Restriction of rows]
\label{lem:restriction-rows}
Let \(N\lhd G\), and fix a block \(b\) of \(N\).
Suppose that every block of \(G\) covering \(b\) has at most
\(L\) simple modules and decomposition numbers at most \(c\).
Suppose also that every simple module in these blocks restricts to
\(N\) with constituent multiplicities at most \(u\).
Then \(b\) has decomposition numbers at most
\(Lcu\).
If \(G/N\) is cyclic, one may take \(u=1\).
If \(C\le Z(G)\), one may instead take \(u=[G:NC]\).
\end{lemma}

\begin{proof}
Choose an ordinary irreducible character upstairs whose restriction
contains a given ordinary irreducible \(\chi\) downstairs.
Decomposition commutes with restriction.  On a fixed Brauer
constituent \(\varphi\) downstairs, the restriction of the
decomposition upstairs has multiplicity at most \(Lcu\).
Its other expression contains \(\chi^0\) with positive integral
multiplicity.  Nonnegativity therefore bounds every
\(d_{\chi\varphi}\) by \(Lcu\).

For the cyclic assertion, a simple \(N\)-type extends to its inertia
group when the quotient is cyclic.  Choose an intertwiner for a
cyclic generator; its power differs from the required action by a
scalar.  Rescale it using a root of that scalar in \(k\).
The remaining multiplicity space is a simple module for a cyclic
quotient, hence is one-dimensional even when \(p\) divides its order.
Clifford theory now gives multiplicity-free restriction.

For the central-factor assertion, let \(X\) be a simple
\(G\)-module and \(S\) a constituent of its restriction to \(N\).
The scalar character by which \(C\) acts on \(X\) extends \(S\)
to \(NC\): its restriction to \(N\cap C\) agrees with the
action on \(S\).  This extension is invariant under the inertia
group \(I\) of \(S\).  The Clifford multiplicity space is
therefore simple for a twisted group algebra of \(I/NC\).
Its dimension is at most \([I:NC]\le [G:NC]\), giving the
claimed restriction bound.
\end{proof}
